\startsectionlevel[title={{\bf 87. Japanese Structure INVERTED: strange life of しか}},reference={japanese-structure-inverted-strange-life-of-しか}]

\goto{{\bf Japanese Structure INVERTED: the strange life of しか. How it really works. Lesson 87}}[url(https://www.youtube.com/watch?v=iEnUH0L6VYs&list=PLg9uYxuZf8x_A-vcqqyOFZu06WlhnypWj&index=91&ab_channel=OrganicJapanesewithCureDolly)]

こんにちは。

Today, I'm afraid I don't have any reindeer for you, but I do have a few deer. We're going to talk about the word シカ / 鹿, which, as you probably know, means deer, ::: info Just in case, this しか is different to the particle しか of this lesson, hence the Katakana, there is a lot of different しか nouns and whatnot with their own Kanji, but only one particle しか ::: {\externalfigure[../media/image754.webp]}

{\bf but there's another しか element in Japanese which is a particle} {\bf and which has a rather unusual effect on the structure of sentences.}

So we're going to look into that effect. One of my commenters reminded me of a rather charming Japanese word play:

ならならしかしかしかられない, / {\em 奈良なら鹿しか叱られない (how it might be in Kanji form)}

{\externalfigure[../media/image140.webp]}

and in natural English we would translate this as {\bf In Nara, only the deer get scolded.}

It's useful to know here that {\bf Nara in the Kansai region is famous for its deer.} People go to Nara to see the deer (and for other things --- it's a lovely historical city).

{\bf Strictly, 奈良なら doesn't mean in Nara.} {\bf It means if it's Nara or, more naturally, in the case of Nara.} But what does the rest of the structure mean?

This is what my commenter was really asking. {\bf The question was, “What is the A-car of this sentence?} {\bf Does it have a zero-A-car somewhere, or does it even have an A-car at all?“}

{\externalfigure[../media/image341.webp]}

What is the A-car of a sentence like this? {\bf Now, the problem is partly that しか has} {\bf an unusual effect on the sentence structure,} {\bf but it's also that the receptive helper is used here} which makes it look a bit more complicated than it is.

\startsectionlevel[title={しか},reference={しか}]

So let's start off with a very simple しか sentence. さくら{\bf しか}いない。

In natural English we'd translate this as, There's nobody here but Sakura.

{\externalfigure[../media/image774.webp]}

{\bf So what's しか doing here?} {\bf しか does two things in a sentence like this.} {\bf It knocks out the が-particle just like か,} so this isn't something that we're not already familiar with.

{\bf When you have か in a sentence where there would also be a が,} {\bf we never say がか or かが.} {\bf か displaces が, but the が is still logically there.} {\bf That's what しか also does.}

\thinrule

{\bf But the other thing しか does is a little bit more unusual.}

{\externalfigure[../media/image591.webp]}

I don't know if you've ever used Photoshop, but in Photoshop you can invert a selection. {\bf What happens is that you select an object,}

{\externalfigure[../media/image626.webp]}

{\bf you press the key selection to invert that selection} {\bf and what happens is that everything outside of that object is selected.}

{\externalfigure[../media/image1078.webp]}

{\bf The object is the only thing there now that isn't selected.}

\thinrule

{\bf This is exactly what しか does.} {\bf While が, you might say, selects an object, a noun,} {\bf and marks it as the A-car, the subject of the sentence,} {\bf しか selects that noun, inverts the selection} {\bf and marks that as the A-car of the sentence.}

{\bf So everything other than that selected noun} {\bf is now the A-car of the sentence.}

So, さくら{\bf しか}いない literally means {\bf Everyone other than} Sakura is not here.

In the deer sentence, which looks a little more complicated, it's exactly the same thing: 奈良なら鹿{\bf しか}叱られない (if it's Nara, {\bf everything other than} the deer are not scolded).

{\externalfigure[../media/image215.webp]}

{\bf So again, we're selecting the deer,} {\bf inverting the selection to everything other than the deer,} {\bf and then that becomes the A-car of the sentence, which is not scolded.}

\thinrule

Now, {\bf there's also an implicit relevance clause in this}, so when we say さくらしかいない {\bf we don't mean there's nothing here but Sakura:} no trees, no bushes, no flying saucers.

{\bf We mean there's nobody here but Sakura.}

{\externalfigure[../media/image1034.webp]}

{\bf Now, the いない partly tells us that}, but it isn't just that, because {\bf it also }doesn't mean** that there are no bunny rabbits here{\bf , }there are no birds, there are no capybara** (did I pronounce that right?)

{\externalfigure[../media/image1029.webp]}

It should simply select the relevance that {\bf no other Sakura / no other people} are there.

{\bf So しか inverts the selection, but it's also} {\bf what you might call a smart invert: it selects for relevance.}

\stopsectionlevel

\startsectionlevel[title={しか in sentences with が-marked subject},reference={しか-in-sentences-with-が-marked-subject}]

{\bf However, in some しか sentences there is actually} {\bf a が-marked subject, so what's going on here?}

{\externalfigure[../media/image961.webp]}

Let's take a look at one. Suppose we say, タオル{\bf が}一枚{\bf しか}ありません.

Again, in natural English: There's {\bf only} one towel here. ::: info Notice how in Japanese they use ありません after しか here. ::: {\bf The が-marked subject is towel or towels.} {\bf In Japanese, of course, there's no distinction between those two things.}

And as I explained in my lesson on counters\goto{{[}71{]}}[url(./71-japanese-counters-3-simple-rules.md)], {\bf a counter in this kind of sentence is working as an adverb.} {\bf It's telling us more about the engine of the sentence, the verb.}

So if we say 泥棒が{\bf 三人}いる, we're saying robbers{\em (=subject)} exist three-{\bf person-ly}. If we say タオルが{\bf 一枚}ある, we're saying towel{\em (=subject)} {\bf one-flat-thing-ly} exists.

If we say タオルが{\bf 二枚}ある, we're saying towel{\em (=subject)} {\bf two-flat-thing-ly} exists.

So when we say タオルが{\bf 一枚}しかありません, {\bf the subject is towels}, {\bf then we have the counter 一枚 and that is marked by しか.} {\bf So we're selecting one flat thing, one towel,} {\bf and then inverting the selection to all towels.}

\thinrule

{\bf Not all flat things, because we've already marked the subject as towel.} {\bf So we're inverting the counter, which is working adverbially,} {\bf from one towel to every towel but that one towel.}

So we're saying towels, {\bf everything but one}, non-exists. In natural English, There's {\bf only one} towel.

{\em タオルが{\bf 一枚しか}ありません}

And that's how what we might call the inverse-selection structure of a しか sentence works.

\stopsectionlevel

\startsectionlevel[title={しかない (Colloquial)},reference={しかない-colloquial}]

{\bf And I should just add at the end that} {\bf there is a use of しか that's more colloquial.} {\bf And this is when we say しかない.}

{\externalfigure[../media/image25.webp]}

Now, in an ordinary しか sentence obviously we just say この古い車{\bf しかない}, which means in natural English There's only this old car.

{\bf We're selecting the car, inverting the selection:} {\bf all cars other than} this old car {\bf non-exist}.

\thinrule

{\bf But we can also put that しかない at the end of a logical clause} {\bf or a verb standing in for a logical clause.} {\bf Now, this isn't strictly grammatical, but colloquially it often happens.}

So if you hear in an anime somebody saying 逃げる{\bf しかない!}, they're saying (we) {\bf gotta} run! in natural English. What they're literally saying is, {\bf we take this 逃げる,} {\bf we mark it with しか so that what we're selecting} {\bf is not 逃げる (run) but everything else other than run.}

And again, {\bf the relevancy smart selection comes in here.} What we're really saying is {\bf every course of action other than} running {\em (away)} {\bf doesn't exist}.

Well, of course it does exist, {\bf this is very colloquial}, but as far as we're concerned it doesn't exist: {\bf there's nothing for it but to run} {\em (away)}.

\thinrule

{\bf In English} we might say, {\bf It's run or nothing}, and again, {\bf that's ungrammatical and for the same reason} {\bf that the Japanese is ungrammatical, that run isn't a noun.}

{\bf But, when that thing's coming after you, who cares about grammar?.} ::: info Some useful comments I guess under the video\ldots{}

{\externalfigure[../media/image124.webp]}

{\externalfigure[../media/image477.webp]}

{\em One miscellaneous about くせに (癖に?), for the full lesson, she made \goto{{\bf this video}}[url(https://www.youtube.com/watch?v=QvuNXIYqFNM&pp=ugMICgJqYRABGAHKBRRjdXJlIGRvbGx5IOOBj-OBm-OBqw\%3D\%3D)]} {\em / Lesson 93}{\externalfigure[../media/image556.webp]} :::

\stopsectionlevel

\stopsectionlevel
